% ===== 原文 pp.152–154（PDF p.162–164） =====
% [原文 p.152]

\begin{proposition}[7] 设 $(M,G)$ 是带 Levi-Civita 联络 $\nabla$ 的（伪）黎曼流形。设 $Y\subseteq M$ 是配备其世界叶上闭 2-形式 $F$ 的子流形，并设 $\cR_+$ 是由式 (9.2) 定义的子丛。若 $\cR_+$ 可积，则对一切 $u\in\Gamma(TY)$ 与 $s\in\Gamma(\cR_+)$ 有
\[
  (e^*\nabla)_u s\in\Gamma(\cR_+).
\]
这里 $e^*\nabla$ 是 $\nabla$ 沿嵌入映射 $e:Y\to M$ 的拉回。用文字说，子丛 $\cR_+$ 关于 $TY$ 是 $\nabla$-平行的。\end{proposition}

% [原文 p.153]
\begin{proof}
我们将利用 $\cR_+$ 的对合性（它由可积性推出）来建立公式
\begin{equation}
  2G((e^*\nabla)_u s,t)=\dd F(u,s,t) \tag{D.2}
\end{equation}
对所有 $u\in\Gamma(TY)$ 与 $s,t\in\Gamma(\cR_+)$ 成立；这里我们始终把 $TY$ 与 $\cR_+$ 的截面看作拉回丛 $e^*TM$ 的截面，该丛配以拉回联络 $e^*\nabla$。由于按假设 $\dd F=0$，该公式将蕴含
\[
  (e^*\nabla)_u s\in\Gamma(\cR_+^\perp)=\Gamma(\cR_+).
\]
在填满空间的情形，我们已注意到我们的流形是（半）para-Kähler 的，而该公式化为 [48] 中定理 1 里出现的一个方程；证明将循相似的思路进行。

为证明 (D.2)，我们先写下 Kosczul 公式
\[
  2G((e^*\nabla)_u s,t)=u(G(s,t))+s(G(u,t))-t(G(u,s))+G([u,s],t)-G([u,t],s)-G([s,t],u).
\]
由于 $s,t\in\Gamma(\cR_+)$，我们可以把上式右端除最后一项之外的所有项中的 $G$ 换成 $-F$。然而，由于 $\cR_+$ 是对合的，$[s,t]\in\cR_+$，故在最后一项中可以换成 $+F$：
\[
  2G((e^*\nabla)_u s,t)=-u(F(s,t))-s(F(u,t))+t(F(u,s))-F([u,s],t)+F([u,t],s)-F([s,t],u).
\]
在右端我们得到了 $F$ 的外微分的正确形式，只是第一项 $-u(F(s,t))$ 的符号不同。然而注意，$\cR_+$ 关于 $G$ 与 $F$ 都是迷向子丛，故该项为零，于是可以翻转其符号。公式 (D.2) 得证。\end{proof}

% [Tutor Note: p.153 原文把 Koszul 公式印作 “Kosczul formula”，是笔误：这是 Levi-Civita 联络的经典 Koszul 公式（Koszul 为人名，无 c）；同页还有一处标点笔误 “Since $s, t\in\Gamma(\cR_+)$. we can replace $\dots$”（句号应为逗号）。]

% [Tutor Note: p.153 由 (D.2) 到结论的省略推理：$\dd F=0$ 使 (D.2) 的右端为零，故 $G((e^*\nabla)_u s,t)=0$ 对一切 $t\in\Gamma(\cR_+)$ 成立，即 $(e^*\nabla)_u s$ 落在 $\cR_+$ 关于 $G$ 的正交补中；再由命题 6，$\cR_+$ 是 $e^*TM$ 的拉格朗日子丛（中间维数、迷向），故 $(\cR_+)^\perp=\cR_+$，结论成立。]

\begin{proposition}[8] 对一切允许变分，项 $G_{IJ}\psi_+^I\delta^{\mathrm{cov}}\psi_+^J$ 为零。\end{proposition}

\begin{proof}
我们把费米边界条件写成
\[
  R(\phi)^I{}_J\psi_+^J=\psi_+^I,
\]
其中 $R:e^*TM\to e^*TM$ 是依赖于 $\phi$ 的矩阵，其平方等于恒等。

% [原文 p.154]
对上式变分，并用协变变分表示，得到
\[
  (\nabla_KR^I{}_J)\delta\phi^K\psi_+^J+R^I{}_J\delta^{\mathrm{cov}}\psi_+^J=\delta^{\mathrm{cov}}\psi_+^I,
\]
或者，由于对允许变分而言 $\delta\phi$ 是切向量、且 $\psi_+$ 取值于 $\cR_+$，
\[
  (\nabla_KR^I{}_J)\pi^K{}_LP^J{}_M\delta\phi^L\psi_+^M+R^I{}_J\delta^{\mathrm{cov}}\psi_+^J=\delta^{\mathrm{cov}}\psi_+^I,
\]
其中 $\pi$ 是到 $TY$ 上的投影算子，$P=\frac{1}{2}(1+R)$ 是到 $\cR_+$ 上的投影算子。然而由命题 7，第一项为零。于是我们有
\[
  G_{IJ}\psi_+^I\delta^{\mathrm{cov}}\psi_+^J=(G_{IJ}P^I{}_KP^J{}_L)\psi_+^K\delta^{\mathrm{cov}}\psi_+^L=0,
\]
因为 $\cR_+$ 是迷向子丛。\end{proof}

% [Tutor Note: p.154 两处补出。(i) $P=\frac{1}{2}(1+R)$ 确是投影算子：由 $R^2=1$ 得 $P^2=\frac{1}{4}(1+2R+R^2)=\frac{1}{2}(1+R)=P$。(ii) 末式为零的依据：$\psi_+$ 取值于 $\cR_+$ 即 $\psi_+=P\psi_+$，故 $G_{IJ}\psi_+^I\delta^{\mathrm{cov}}\psi_+^J=(P^TGP)_{KL}\psi_+^K\delta^{\mathrm{cov}}\psi_+^L$；而 $\cR_+$ 关于 $G$ 迷向正是 $P^TGP=0$，于是全项为零。]

% [Tutor Note: 附录 D 到此结束（p.154 余下为空白）。原文只给出命题 1、3、6、7、8 的证明，其中命题 1 的证明用式 (D.1)，命题 7 的证明用式 (D.2)，与 tags.txt 中 D.1（p.151）、D.2（p.153）的清单一致。]
